%% ============================================================
%%  Rapport de Projet de Fin d'Études — Tekup
%%  Méthodologie CRISP-DM — Génération 3D par IA
%%  Étudiant : Houcem Belkhiria
%% ============================================================
\documentclass[12pt,a4paper,oneside]{report}

%% ---- Encodage & langue ----
\usepackage[utf8]{inputenc}
\usepackage[T1]{fontenc}
\usepackage[french]{babel}
\usepackage{lmodern}

%% ---- Mise en page ----
\usepackage[top=2.5cm, bottom=2.5cm, left=3cm, right=2.5cm]{geometry}
\usepackage{setspace}
\onehalfspacing
\usepackage{fancyhdr}
\usepackage{emptypage}

%% ---- Couleurs & graphiques ----
\usepackage[dvipsnames,svgnames,x11names]{xcolor}
\definecolor{tekupblue}{RGB}{0,84,166}
\definecolor{crispyellow}{RGB}{255,180,0}
\definecolor{crispgreen}{RGB}{34,139,34}
\definecolor{lightgray}{RGB}{245,245,245}
\definecolor{codegray}{RGB}{240,240,240}
\usepackage{graphicx}
\usepackage{tikz}
\usetikzlibrary{shapes.geometric, arrows.meta, positioning,
                fit, backgrounds, decorations.pathmorphing}
\usepackage{pgfplots}
\pgfplotsset{compat=1.18}

%% ---- Tableaux ----
\usepackage{booktabs}
\usepackage{tabularx}
\usepackage{multirow}
\usepackage{array}
\usepackage{longtable}
\usepackage{colortbl}

%% ---- Code source ----
\usepackage{listings}
\lstset{
  backgroundcolor=\color{codegray},
  basicstyle=\ttfamily\footnotesize,
  breaklines=true,
  captionpos=b,
  commentstyle=\color{OliveGreen},
  keywordstyle=\color{tekupblue}\bfseries,
  stringstyle=\color{BrickRed},
  numberstyle=\tiny\color{gray},
  numbers=left, numbersep=5pt,
  frame=single, rulecolor=\color{gray!30},
  tabsize=2, xleftmargin=10pt,
}

%% ---- Hyperliens ----
\usepackage[colorlinks=true, linkcolor=tekupblue, citecolor=tekupblue,
            urlcolor=tekupblue, bookmarks=true]{hyperref}

%% ---- Bibliographie ----
\usepackage[backend=biber, style=ieee, sorting=none]{biblatex}
\addbibresource{references.bib}

%% ---- Figures & légendes ----
\usepackage{caption}
\usepackage{subcaption}
\usepackage{float}
\captionsetup{font=small, labelfont=bf}

%% ---- Mathématiques ----
\usepackage{amsmath, amssymb}

%% ---- Boîtes & listes ----
\usepackage{enumitem}
\usepackage{mdframed}
\usepackage{tcolorbox}
\tcbuselibrary{skins, breakable}

%% ---- Boîte CRISP-DM phase ----
\newtcolorbox{crispbox}[2][]{
  colback=tekupblue!6, colframe=tekupblue,
  fonttitle=\bfseries\color{white},
  title={#2}, #1,
  before upper={\parindent0pt},
}

%% ---- Titre des chapitres ----
\usepackage{titlesec}
\titleformat{\chapter}[display]
  {\normalfont\huge\bfseries\color{tekupblue}}
  {\chaptertitlename\ \thechapter}{16pt}{\Huge}
\titlespacing*{\chapter}{0pt}{-20pt}{20pt}

%% ---- En-têtes / pieds de page ----
\pagestyle{fancy}
\fancyhf{}
\fancyhead[L]{\small\nouppercase{\leftmark}}
\fancyhead[R]{\small Houcem Belkhiria — Tekup 2025}
\fancyfoot[C]{\thepage}
\renewcommand{\headrulewidth}{0.4pt}

%% ================================================================
\begin{document}

%% ----------------------------------------------------------------
%%  PAGE DE GARDE
%% ----------------------------------------------------------------
\begin{titlepage}
\centering

{\large \textbf{République Tunisienne}\\
Ministère de l'Enseignement Supérieur et de la Recherche Scientifique}\\[0.4cm]

\begin{tikzpicture}
  \draw[tekupblue, line width=2pt] (0,0) rectangle (14,2.2);
  \node[text=tekupblue, font=\bfseries\large] at (7,1.5)
        {Université Tekup};
  \node[font=\normalsize] at (7,0.6)
        {Institut Supérieur des Technologies de l'Information};
\end{tikzpicture}

\vspace{0.6cm}
\rule{\linewidth}{1pt}\\[0.2cm]
{\LARGE \textbf{Projet de Fin d'Études}}\\[0.1cm]
{\large Licence Nationale en Informatique}\\[0.1cm]
\rule{\linewidth}{1pt}

\vspace{0.5cm}

{\color{tekupblue}\Huge\bfseries
Conception et Développement d'un\\[0.2cm]
Microservice d'IA pour la Génération\\[0.2cm]
Automatisée d'Assets 3D Unity via MCP
}

\vspace{0.4cm}
\begin{tikzpicture}
  \node[draw=crispyellow, fill=crispyellow!20, rounded corners=5pt,
        inner sep=8pt, font=\large\bfseries] {
    Méthodologie CRISP-DM
  };
\end{tikzpicture}

\vspace{0.6cm}
\begin{mdframed}[linecolor=tekupblue, linewidth=1pt,
                  leftmargin=30pt, rightmargin=30pt]
\centering
\textbf{Réalisé par :} Houcem Belkhiria\\[0.3cm]
\textbf{Encadrant académique :} \textit{(nom de l'encadrant)}\\[0.1cm]
\textbf{Encadrant professionnel :} \textit{(nom de l'encadrant)}
\end{mdframed}

\vspace{0.6cm}
\begin{tikzpicture}
  \node[draw=tekupblue, rounded corners=6pt, inner sep=10pt,
        fill=tekupblue!10] {
    \begin{tabular}{ll}
      \textbf{Spécialité :}           & Génie Logiciel \& Intelligence Artificielle \\
      \textbf{Année universitaire :}  & 2024 -- 2025 \\
      \textbf{Période de stage :}    & Février 2025 -- Juillet 2025 \\
      \textbf{Date de soutenance :}   & Début Juillet 2025 \\
    \end{tabular}
  };
\end{tikzpicture}

\vfill
{\small \textit{Soutenu devant le jury composé de :}\\
Mme/M. \underline{\hspace{4cm}} \quad Président(e)\\
Mme/M. \underline{\hspace{4cm}} \quad Examinateur/trice\\
Mme/M. \underline{\hspace{4cm}} \quad Encadrant(e)}
\end{titlepage}

%% ----------------------------------------------------------------
%%  DÉDICACE
%% ----------------------------------------------------------------
\cleardoublepage\thispagestyle{empty}
\vspace*{6cm}
\begin{center}
\textit{\Large Dédicace}\\[1cm]\large
\textit{À mes parents, source inépuisable de soutien et d'amour,\\
dont les sacrifices ont rendu ce parcours possible.}\\[0.6cm]
\textit{À mes amis et camarades de Tekup,\\
qui ont partagé avec moi chaque doute et chaque victoire.}\\[0.6cm]
\textit{À tous ceux qui croient que la technologie\\
peut rendre la création accessible à tous.}
\end{center}

%% ----------------------------------------------------------------
%%  REMERCIEMENTS
%% ----------------------------------------------------------------
\cleardoublepage
\chapter*{Remerciements}
\addcontentsline{toc}{chapter}{Remerciements}

Je tiens à remercier chaleureusement mon encadrant académique pour sa
disponibilité, ses conseils avisés et l'autonomie qu'il m'a accordée
tout au long de ce projet.

Mes remerciements vont également à l'équipe pédagogique de Tekup pour la
qualité de la formation en génie logiciel et intelligence artificielle
dispensée au cours de ces trois années.

Enfin, je remercie ma famille dont le soutien moral indéfectible m'a permis
de mener ce projet ambitieux à son terme.

%% ----------------------------------------------------------------
%%  RÉSUMÉ
%% ----------------------------------------------------------------
\cleardoublepage
\chapter*{Résumé}
\addcontentsline{toc}{chapter}{Résumé}

Ce projet de fin d'études présente la conception et le développement d'un
\textbf{microservice d'intelligence artificielle} capable de transformer des
données non structurées (PDF, emails, images) en objets 3D, puis de les
injecter automatiquement dans l'éditeur Unity via le protocole \textbf{MCP}
(\textit{Model Control Protocol}).

La méthodologie \textbf{CRISP-DM} (\textit{Cross-Industry Standard Process for
Data Mining}) structure le projet en six phases : compréhension du domaine,
compréhension des données, préparation des données, modélisation, évaluation
et déploiement.

Techniquement, le système repose sur un backend \textbf{FastAPI/Celery/Redis},
les pipelines de génération 3D \textbf{Hunyuan3D-2} (diffusion DiT + VAE 3D),
un cache vectoriel \textbf{ChromaDB}, et un frontend \textbf{React 19}. Le
pont Unity est assuré par le script éditeur \texttt{SpawnBridge.cs} couplé
à la bibliothèque \textbf{glTFast 6.9.0}.

\textbf{Mots-clés :} CRISP-DM, génération 3D, IA générative, Hunyuan3D,
FastAPI, MCP, Unity, ChromaDB, microservice.

\vspace{0.8cm}
\noindent\rule{\linewidth}{0.4pt}
\begin{center}\textbf{Abstract}\end{center}

This final-year project presents the design and development of an
\textbf{AI microservice} that transforms unstructured data (PDFs, emails,
images) into 3D objects and automatically injects them into the Unity Editor
via the \textbf{Model Control Protocol (MCP)}. The \textbf{CRISP-DM}
methodology structures the project into six phases. Technically, the system
relies on a \textbf{FastAPI/Celery/Redis} backend, \textbf{Hunyuan3D-2}
diffusion pipelines, a \textbf{ChromaDB} vector cache, and a \textbf{React 19}
frontend. The Unity bridge is provided by the \texttt{SpawnBridge.cs} editor
script coupled with \textbf{glTFast}.

\textbf{Keywords:} CRISP-DM, 3D generation, generative AI, Hunyuan3D, FastAPI,
MCP, Unity, ChromaDB, microservice.

%% ----------------------------------------------------------------
%%  TABLES
%% ----------------------------------------------------------------
\cleardoublepage\tableofcontents
\cleardoublepage\listoffigures
\addcontentsline{toc}{chapter}{Table des figures}
\cleardoublepage\listoftables
\addcontentsline{toc}{chapter}{Liste des tableaux}

%% ----------------------------------------------------------------
%%  ABRÉVIATIONS
%% ----------------------------------------------------------------
\cleardoublepage
\chapter*{Liste des abréviations}
\addcontentsline{toc}{chapter}{Liste des abréviations}
\begin{tabular}{@{}lp{10.5cm}@{}}
\toprule
\textbf{Sigle} & \textbf{Signification} \\\midrule
API     & Application Programming Interface \\
CRISP-DM& Cross-Industry Standard Process for Data Mining \\
DiT     & Diffusion Transformer \\
GLB     & GL Transmission Format Binary \\
GPU     & Graphics Processing Unit \\
HTTP    & HyperText Transfer Protocol \\
JSON    & JavaScript Object Notation \\
LLM     & Large Language Model \\
MCP     & Model Control Protocol \\
MPS     & Metal Performance Shaders (Apple Silicon) \\
PBR     & Physically Based Rendering \\
PFE     & Projet de Fin d'Études \\
REST    & Representational State Transfer \\
SPA     & Single Page Application \\
T2I     & Text-to-Image \\
UUID    & Universally Unique Identifier \\
US      & User Story \\
VRAM    & Video Random Access Memory \\
\bottomrule
\end{tabular}

%% ================================================================
%%  INTRODUCTION GÉNÉRALE
%% ================================================================
\cleardoublepage
\chapter*{Introduction générale}
\addcontentsline{toc}{chapter}{Introduction générale}
\markboth{Introduction générale}{}

La modélisation 3D est au cœur du développement de jeux vidéo, de la
simulation industrielle et de la réalité virtuelle. Pourtant, produire un
asset 3D de qualité reste long, coûteux et réservé aux artistes spécialisés.
L'émergence des modèles de diffusion 3D — en particulier \textbf{Hunyuan3D-2}
de Tencent — ouvre la voie à une génération automatisée depuis une simple
image ou une phrase descriptive.

Ce projet de fin d'études a pour ambition de \textbf{fermer la boucle} entre
la génération IA et l'utilisation directe dans Unity : l'utilisateur décrit
ou photographie un objet, la plateforme génère le maillage 3D, et l'objet
apparaît automatiquement dans la scène Unity — sans aucune manipulation manuelle.

\subsection*{Problématique}
\textit{Comment concevoir un microservice d'IA agentique capable de
transformer des données non structurées en assets 3D utilisables dans Unity,
en garantissant des performances acceptables sur matériel grand public et
une expérience utilisateur fluide ?}

\subsection*{Méthodologie adoptée}
La méthodologie \textbf{CRISP-DM} (\textit{Cross-Industry Standard Process
for Data Mining}), standard industriel pour les projets IA/Data, structure
ce travail en six phases itératives :

\begin{figure}[H]
\centering
\begin{tikzpicture}[
  phase/.style={draw=tekupblue, fill=tekupblue!10, rounded corners=5pt,
                text width=2.8cm, align=center, minimum height=1cm,
                font=\small\bfseries},
  arr/.style={-Stealth, tekupblue, thick}
  ]
  \node[phase, fill=crispyellow!30, draw=crispyellow!80!black] (p1) at (0,0)
        {1. Compréhension\\du domaine};
  \node[phase] (p2) at (4,0)   {2. Compréhension\\des données};
  \node[phase] (p3) at (8,0)   {3. Préparation\\des données};
  \node[phase] (p4) at (8,-2.5){4. Modélisation};
  \node[phase] (p5) at (4,-2.5){5. Évaluation};
  \node[phase, fill=crispgreen!20, draw=crispgreen!60!black] (p6) at (0,-2.5)
        {6. Déploiement};

  \draw[arr] (p1) -- (p2);
  \draw[arr] (p2) -- (p3);
  \draw[arr] (p3) -- (p4);
  \draw[arr] (p4) -- (p5);
  \draw[arr] (p5) -- (p6);
  \draw[arr, dashed] (p5) to[bend left=20] (p1);
  \node[font=\scriptsize, color=gray] at (4,-1.3) {itération};
\end{tikzpicture}
\caption{Cycle CRISP-DM appliqué au projet}
\label{fig:crispdm}
\end{figure}

\subsection*{Structure du rapport}
Chaque chapitre correspond à une phase CRISP-DM, complétée par l'introduction
(contexte), la conclusion et les perspectives.

%% ================================================================
%%  PHASE 1 — COMPRÉHENSION DU DOMAINE
%% ================================================================
\chapter{Phase 1 — Compréhension du domaine}

\begin{crispbox}{CRISP-DM Phase 1 : Business Understanding}
Objectif : définir les objectifs métier, évaluer la situation existante,
fixer les critères de succès et élaborer le plan de projet (backlog).
\end{crispbox}

\section{Contexte et objectifs métier}

\subsection{Contexte industriel}
Le marché mondial de la modélisation 3D est évalué à \textbf{14 milliards USD}
d'ici 2030 (Grand View Research). Unity Technologies recense plus de
\textbf{1,5 million} de développeurs actifs chaque mois. La demande en assets
3D prêts à l'emploi est massive mais la création manuelle représente un
goulot d'étranglement majeur.

\subsection{Objectifs du microservice}
L'objectif principal est de développer une architecture microservice autonome
et \textbf{agentique} capable de :
\begin{enumerate}
  \item \textbf{Ingérer} des fichiers multi-formats (PDF, Email, images).
  \item \textbf{Extraire} des métadonnées structurées via un LLM.
  \item \textbf{Générer} un objet 3D (approche neuronale : Image/Text-to-3D ;
        approche procédurale : script C\#).
  \item \textbf{Piloter Unity via MCP} : le microservice agit comme un client
        MCP qui ordonne à l'éditeur Unity (serveur MCP) d'instancier et de
        placer l'objet dans la scène active — sans clic utilisateur.
\end{enumerate}

\subsection{Critères de succès}
\begin{table}[H]
\centering
\caption{Critères de succès CRISP-DM}
\label{tab:success}
\begin{tabularx}{\linewidth}{@{}lXl@{}}
\toprule
\textbf{Critère} & \textbf{Description} & \textbf{Seuil} \\
\midrule
Performance & Temps de génération GLB sans texture & $\leq$ 2 min \\
Qualité géom. & Absence de floaters et de faces dégénérées & $\geq$ 80\% \\
Intégration & Spawn Unity sans erreur & 100\% \\
Disponibilité & Uptime de l'API & $\geq$ 99\% \\
UX & Génération depuis l'interface sans CLI & Oui \\
\bottomrule
\end{tabularx}
\end{table}

\section{Étude de l'existant}

\subsection{Solutions de génération 3D}
\begin{longtable}{@{}p{2.8cm}p{3.6cm}p{3.5cm}p{4.2cm}@{}}
\caption{Comparatif des solutions de génération 3D}\label{tab:existant}\\
\toprule
\textbf{Outil} & \textbf{Approche} & \textbf{Avantages} & \textbf{Limites} \\
\midrule\endfirsthead
\toprule\textbf{Outil} & \textbf{Approche} & \textbf{Avantages} & \textbf{Limites} \\
\midrule\endhead
TripoSR      & NeRF/Transformer   & Rapide ($<$1 s)      & Maillage bas détail \\
Shap·E       & Diffusion implicite & Open-source         & Qualité insuffisante \\
Meshy.ai     & API cloud          & Simple d'usage       & Payant, latence réseau \\
Zero123++    & Diffusion multi-vues& Bonne précision     & Nécessite GPU puissant \\
\textbf{Hunyuan3D-2} & DiT + VAE 3D & Haute qualité, texture & 8 GB RAM minimum \\
\bottomrule
\end{longtable}

\subsection{Le Model Context Protocol (MCP)}
Le \textbf{MCP} est un standard ouvert proposé par Anthropic en 2024 pour
standardiser la connexion entre assistants IA et outils locaux. Dans notre
architecture :
\begin{itemize}
  \item \textbf{Microservice Python} = Client MCP — appelle des outils.
  \item \textbf{Unity Editor} = Serveur MCP — expose des outils C\# :
        \texttt{SpawnGLB(url, position)}, \texttt{SetTransform(...)},
        \texttt{LogMessage(text)}.
\end{itemize}
Ce protocole remplace le \textit{polling} de fichiers par une communication
bidirectionnelle robuste, transformant le projet d'un simple générateur en
\textbf{assistant intelligent agentique}.

\section{Backlog Agile — 5 mois / 8 sprints}

\subsection{User Stories par phase}

\begin{table}[H]
\centering
\caption{Backlog complet du projet}
\label{tab:backlog}
\begin{tabularx}{\linewidth}{@{}llXl@{}}
\toprule
\rowcolor{tekupblue!15}
\textbf{ID} & \textbf{Phase} & \textbf{User Story} & \textbf{Sprint} \\
\midrule
\rowcolor{gray!8}
US-1.1 & Infrastructure & Setup Docker-Compose (FastAPI, Redis, Celery) & 1 \\
US-1.2 & Ingestion      & Pipeline \texttt{unstructured} pour PDF et emails & 1 \\
US-1.3 & Ingestion      & Filtres Regex : suppression signatures / bruit & 1 \\
US-1.4 & API            & Endpoint \texttt{/upload} asynchrone + \texttt{task\_id} & 2 \\
\midrule
\rowcolor{tekupblue!8}
US-2.1 & LLM            & Intégration Llama-3-8B-Instruct (llama-cpp) & 2 \\
US-2.2 & LLM            & Prompt système pour sortie JSON forcée (Schema Enforcement) & 2 \\
US-2.3 & LLM            & Extraction de tableaux de dimensions depuis PDF & 3 \\
US-2.4 & Validation     & Validation Pydantic des métadonnées extraites & 3 \\
\midrule
\rowcolor{gray!8}
US-3.1 & 3D Neuronal    & Intégration Hunyuan3D-2 Image-to-3D & 3 \\
US-3.2 & 3D Procédural  & Pipeline scripts C\# via Qwen 2.5 Coder & 4 \\
US-3.3 & 3D             & Conversion et optimisation assets (OBJ → GLB) & 4 \\
US-3.4 & Stockage       & Gestion du stockage et nettoyage des fichiers & 4 \\
\midrule
\rowcolor{tekupblue!8}
US-4.1 & MCP Python     & Client MCP dans le microservice (SDK Anthropic) & 5 \\
US-4.2 & MCP Unity      & Serveur MCP en C\# dans Unity (WebSocket/Stdio) & 5 \\
US-4.3 & MCP Unity      & Définition des outils Unity : SpawnGLB, SetTransform & 5 \\
US-4.4 & Orchestration  & Connexion pipeline Celery → appel outil MCP & 6 \\
\midrule
\rowcolor{gray!8}
US-5.1 & UX Unity       & Fenêtre « Moniteur IA » dans Unity (logs MCP) & 6 \\
US-5.2 & Test E2E       & Test bout en bout : PDF → extraction → 3D → Unity & 6 \\
US-5.3 & Doc            & Rédaction mémoire et documentation architecture & 7 \\
US-5.4 & Optimisation   & Nettoyage code et optimisation VRAM & 8 \\
\bottomrule
\end{tabularx}
\end{table}

\subsection{Planning des sprints}

\begin{table}[H]
\centering
\caption{Planning des 8 sprints --- Février à Juillet 2025 (22 semaines)}
\label{tab:sprints}
\begin{tabularx}{\linewidth}{@{}cllXl@{}}
\toprule
\textbf{Spr.} & \textbf{Durée} & \textbf{Période (2025)}
              & \textbf{Objectif principal} & \textbf{US} \\
\midrule
1 & 2 sem. & 03 Fév -- 14 Fév & Infrastructure Docker/FastAPI/Celery + ingestion & US-1.1--1.3 \\
2 & 2 sem. & 17 Fév -- 28 Fév & Upload async + intégration LLM Llama-3 + JSON & US-1.4, 2.1--2.2 \\
3 & 3 sem. & 03 Mar -- 21 Mar & Extraction complexe, Pydantic, Hunyuan3D I2-3D & US-2.3--2.4, 3.1 \\
4 & 3 sem. & 24 Mar -- 11 Avr & Pipeline C\# procédural, conversion GLB, stockage & US-3.2--3.4 \\
5 & 3 sem. & 14 Avr -- 02 Mai & Client MCP Python + Serveur MCP Unity + outils & US-4.1--4.3 \\
6 & 3 sem. & 05 Mai -- 23 Mai & Orchestration Celery→MCP + Moniteur IA + E2E & US-4.4, 5.1--5.2 \\
7 & 3 sem. & 26 Mai -- 13 Jun & Rédaction mémoire + documentation architecture & US-5.3 \\
8 & 3 sem. & 16 Jun -- 04 Jul & Optimisation VRAM + polissage + soutenance & US-5.4 \\
\bottomrule
\end{tabularx}
\end{table}

%% ================================================================
%%  PHASE 2 — COMPRÉHENSION DES DONNÉES
%% ================================================================
\chapter{Phase 2 — Compréhension des données}

\begin{crispbox}{CRISP-DM Phase 2 : Data Understanding}
Objectif : collecter, décrire, explorer et vérifier la qualité des données
d'entrée qui alimenteront les modèles IA.
\end{crispbox}

\section{Sources de données}

Le système ingère plusieurs types de données selon le mode d'utilisation :

\begin{table}[H]
\centering
\caption{Sources de données du système}
\label{tab:data-sources}
\begin{tabularx}{\linewidth}{@{}llXl@{}}
\toprule
\textbf{Source} & \textbf{Format} & \textbf{Description} & \textbf{Usage} \\
\midrule
Image utilisateur & PNG, JPEG & Photo ou rendu d'un objet réel & Image-to-3D \\
Texte libre       & UTF-8     & Description en langage naturel   & Text-to-3D \\
Multi-vues        & PNG, JPEG & 2 à 4 vues (face/dos/gauche/droite) & MultiView-to-3D \\
PDF technique     & PDF       & Cahiers des charges, fiches techniques & LLM + extraction \\
Email             & EML       & Demandes clients avec spécifications & LLM + extraction \\
\bottomrule
\end{tabularx}
\end{table}

\section{Exploration des données d'images}

\subsection{Caractéristiques requises}
Pour que le pipeline de génération 3D produise un résultat de qualité, les
images d'entrée doivent respecter les contraintes suivantes :

\begin{itemize}
  \item \textbf{Résolution} : $\geq 512 \times 512$ pixels (recommandé :
        $1024 \times 1024$).
  \item \textbf{Sujet centré} : l'objet doit occuper au moins 60\% de la
        surface de l'image.
  \item \textbf{Fond contrasté} : fond blanc ou uni pour faciliter la
        suppression d'arrière-plan.
  \item \textbf{Éclairage homogène} : éviter les ombres portées dures.
\end{itemize}

\subsection{Jeu de données de test}
Un corpus de \textbf{20 images de test} a été constitué, couvrant :
véhicules (5), mobilier (5), personnages (4), végétation (3), objets
du quotidien (3).

\begin{table}[H]
\centering
\caption{Distribution du corpus de test}
\label{tab:corpus}
\begin{tabular}{@{}lccl@{}}
\toprule
\textbf{Catégorie} & \textbf{Nb. images} & \textbf{\%} & \textbf{Difficulté} \\
\midrule
Véhicules    & 5 & 25\% & Moyenne \\
Mobilier     & 5 & 25\% & Faible \\
Personnages  & 4 & 20\% & Élevée \\
Végétation   & 3 & 15\% & Élevée \\
Divers       & 3 & 15\% & Variable \\
\bottomrule
\end{tabular}
\end{table}

\section{Exploration des données textuelles}

Les données textuelles (PDF, emails) présentent des caractéristiques
hétérogènes : format tabulaire, mise en page multi-colonnes, présence de
signatures et de clauses légales. La bibliothèque \texttt{unstructured}
permet un partitionnement sémantique pour isoler les sections pertinentes.

\subsection{Qualité des données — problèmes identifiés}

\begin{table}[H]
\centering
\caption{Problèmes de qualité des données textuelles}
\label{tab:dq}
\begin{tabularx}{\linewidth}{@{}lXl@{}}
\toprule
\textbf{Problème} & \textbf{Description} & \textbf{Fréquence} \\
\midrule
Bruit textuel & Signatures, mentions légales, pieds de page & Élevée \\
Tableaux complexes & Dimensions en format tabulaire (L×l×h) & Moyenne \\
Encodage mixte & Fichiers latin-1 ou UTF-16 mal déclarés & Faible \\
Ambiguïté unitaire & mm vs cm vs m non précisés & Moyenne \\
\bottomrule
\end{tabularx}
\end{table}

%% ================================================================
%%  PHASE 3 — PRÉPARATION DES DONNÉES
%% ================================================================
\chapter{Phase 3 — Préparation des données}

\begin{crispbox}{CRISP-DM Phase 3 : Data Preparation}
Objectif : transformer les données brutes en entrées exploitables par les
modèles IA — suppression du bruit, normalisation, encodage, augmentation.
\end{crispbox}

\section{Pipeline de préparation des images}

\subsection{Suppression d'arrière-plan (rembg)}
La première étape de traitement est la suppression automatique de
l'arrière-plan via la bibliothèque \texttt{rembg}, basée sur le réseau
de segmentation \textbf{U2-Net} \cite{rembg}. L'image résultante est une
image RGBA avec fond transparent.

\begin{lstlisting}[language=Python,
  caption={Suppression d'arrière-plan dans hunyuan3d\_service.py}]
from hy3dgen.rembg import BackgroundRemover

self.rembg = BackgroundRemover()

def _remove_background(self, image: Image.Image) -> Image.Image:
    # Conversion RGBA obligatoire avant rembg
    if image.mode != 'RGBA':
        image = image.convert('RGBA')
    return self.rembg(image)
\end{lstlisting}

\subsection{Normalisation et encodage Base64}
Les images sont transmises entre le frontend et le backend en encodage
\textbf{Base64} pour éviter la gestion de fichiers temporaires dans
les requêtes HTTP multipart :

\begin{equation}
\text{image\_b64} = \texttt{base64.b64encode(file.read()).decode('utf-8')}
\label{eq:b64}
\end{equation}

\subsection{Prétraitement multi-vues}
Pour le mode Multi-View-to-3D, chaque vue (face, dos, gauche, droite) est
traitée indépendamment puis redimensionnée à $512\times512$ par le
\texttt{ImageProcessorV2} d'hy3dgen, qui applique également un recadrage
centré avec un ratio de bordure de 0,15.

\section{Pipeline de préparation des documents}

\subsection{Extraction de texte structuré}
La bibliothèque \texttt{unstructured} partitionne les PDF et emails en
éléments sémantiques (\textit{Title}, \textit{NarrativeText},
\textit{Table}, etc.) :

\begin{lstlisting}[language=Python,
  caption={Extraction de documents dans document\_parser.py}]
from unstructured.partition.auto import partition

elements = partition(filename=filepath)
sections = [e.text for e in elements
            if e.category in ("NarrativeText", "Title", "Table")]
\end{lstlisting}

\subsection{Nettoyage et filtrage}
Des filtres Regex suppriment les données non pertinentes :

\begin{lstlisting}[language=Python,
  caption={Filtres de nettoyage du texte}]
import re

NOISE_PATTERNS = [
    r'Confidentiel.*',
    r'Page \d+ sur \d+',
    r'[A-Za-z0-9._%+-]+@[A-Za-z0-9.-]+\.[A-Z|a-z]{2,}',  # emails
    r'Tel\s*:.*',
]

def clean_text(text: str) -> str:
    for pat in NOISE_PATTERNS:
        text = re.sub(pat, '', text, flags=re.IGNORECASE)
    return text.strip()
\end{lstlisting}

\subsection{Schema Enforcement via LLM}
Le LLM (\textbf{Llama-3-8B-Instruct}) est guidé par un prompt système strict
qui force une sortie JSON conforme au schéma Pydantic \texttt{UnityMetadata} :

\begin{lstlisting}[language=Python, caption={Prompt système LLM}]
SYSTEM_PROMPT = """
Tu es un expert en extraction de metadonnees 3D pour Unity.
Reponds UNIQUEMENT avec un objet JSON valide respectant ce schema:
{
  "name": "string",
  "type": "chair|table|vehicle|character|...",
  "color": "string",
  "dimensions": {"width": float, "height": float, "depth": float},
  "material": "string",
  "style": "string"
}
Ne fournis aucune explication. JSON uniquement.
"""
\end{lstlisting}

\subsection{Validation Pydantic}
Chaque sortie JSON est validée par un modèle Pydantic avant d'être transmise
au pipeline de génération :

\begin{lstlisting}[language=Python, caption={Modèle Pydantic UnityMetadata}]
from pydantic import BaseModel, Field
from typing import Optional

class Dimensions(BaseModel):
    width: float = Field(gt=0)
    height: float = Field(gt=0)
    depth: float = Field(gt=0)

class UnityMetadata(BaseModel):
    name: str
    type: str
    color: Optional[str] = None
    dimensions: Optional[Dimensions] = None
    material: Optional[str] = None
    style: Optional[str] = None
\end{lstlisting}

%% ================================================================
%%  PHASE 4 — MODÉLISATION
%% ================================================================
\chapter{Phase 4 — Modélisation}

\begin{crispbox}{CRISP-DM Phase 4 : Modeling}
Objectif : sélectionner et construire les modèles IA, configurer les
hyperparamètres et valider les techniques de génération 3D retenues.
\end{crispbox}

\section{Architecture globale du système}

\begin{figure}[H]
\centering
\begin{tikzpicture}[
  box/.style={draw=tekupblue, rounded corners=4pt, fill=tekupblue!8,
              text width=3cm, align=center, minimum height=1.2cm, font=\small},
  arrow/.style={-Stealth, thick, tekupblue},
  lbl/.style={font=\scriptsize, color=gray}
  ]
  \node[box, fill=green!10, draw=green!60!black] (fe) at (0,0)
        {\textbf{Frontend}\\React 19 / Vite\\Tailwind CSS};
  \node[box] (be) at (4.5,0)
        {\textbf{Backend}\\FastAPI\\Uvicorn};
  \node[box, fill=orange!10, draw=orange!70!black] (wk) at (4.5,-3)
        {\textbf{Worker Celery}\\Hunyuan3D-2\\PyTorch / MPS};
  \node[box, fill=red!10, draw=red!60!black] (rd) at (9,0)
        {\textbf{Redis}\\Broker\\Result store};
  \node[box, fill=purple!10, draw=purple!60!black] (db) at (9,-3)
        {\textbf{ChromaDB}\\Cache vectoriel\\CLIP embeddings};
  \node[box, fill=gray!15, draw=gray!60] (un) at (0,-3)
        {\textbf{Unity Editor}\\SpawnBridge.cs\\glTFast 6.9.0};

  \draw[arrow] (fe) -- node[above, lbl]{HTTP/REST} (be);
  \draw[arrow] (be) -- node[right, lbl]{Celery task} (wk);
  \draw[arrow] (be) -- node[above, lbl]{publish/poll} (rd);
  \draw[arrow] (wk) -- node[right, lbl]{store result} (rd);
  \draw[arrow] (wk) -- node[below, lbl]{CLIP embed.} (db);
  \draw[arrow] (be) -- node[left, lbl]{SpawnRequest JSON\\(protocole MCP)} (un);
\end{tikzpicture}
\caption{Architecture générale — vue déploiement}
\label{fig:arch}
\end{figure}

\section{Modèles IA sélectionnés}

\subsection{Hunyuan3D-2 — Génération de forme (ShapeGen)}

Hunyuan3D-2 \cite{hunyuan3d2024} repose sur une architecture
\textbf{DiT} (\textit{Diffusion Transformer}) opérant sur une représentation
octree multi-résolution. Le VAE 3D encode la géométrie implicite, que le
décodeur de flux marching-cubes convertit en maillage triangulé.

Paramètres clés exposés par l'API :

\begin{table}[H]
\centering
\caption{Hyperparamètres du pipeline ShapeGen}
\label{tab:hyper}
\begin{tabularx}{\linewidth}{@{}llXl@{}}
\toprule
\textbf{Paramètre} & \textbf{Type} & \textbf{Rôle} & \textbf{Défaut} \\
\midrule
\texttt{seed}                  & int   & Reproductibilité stochastique    & 1234 \\
\texttt{num\_inference\_steps} & int   & Étapes de débruitage DiT         & 5 \\
\texttt{guidance\_scale}       & float & Intensité du guidage conditionnel & 5.0 \\
\texttt{octree\_resolution}    & int   & Résolution de la grille 3D       & 128 \\
\texttt{num\_chunks}           & int   & Taille des chunks décodeur       & 8000 \\
\texttt{face\_count}           & int   & Nombre max de faces (réduction)  & 20\,000 \\
\bottomrule
\end{tabularx}
\end{table}

\subsection{Hunyuan3D-2 — Génération de texture (TexGen)}

Le module TexGen génère des textures PBR (\textit{Physically Based Rendering})
par \textbf{inpainting multi-vues} : à partir de l'image de référence, six
vues synthétiques sont générées et projetées sur le maillage par UV mapping.

\subsection{Modèles Text-to-Image (T2I)}

Deux backends T2I sont supportés en fonction des contraintes matérielles :

\begin{table}[H]
\centering
\caption{Backends Text-to-Image disponibles}
\label{tab:t2i}
\begin{tabularx}{\linewidth}{@{}lXll@{}}
\toprule
\textbf{Modèle} & \textbf{Description} & \textbf{VRAM} & \textbf{Langue} \\
\midrule
Hyper-SDXL      & SDXL distillé (4 étapes) & 8 GB & Anglais \\
HunyuanDiT      & DiT bilingue Tencent    & 10 GB & FR + ZH \\
\bottomrule
\end{tabularx}
\end{table}

\subsection{LLM d'extraction — Llama-3-8B-Instruct}

Llama-3 est exécuté localement via \texttt{llama-cpp-python} avec
quantisation \textbf{GGUF Q4\_K\_M} pour réduire l'empreinte mémoire à
$\approx$ 5 GB. Il reçoit le texte extrait et produit un JSON structuré.

\section{Pipeline de génération asynchrone}

\begin{figure}[H]
\centering
\begin{tikzpicture}[
  step/.style={draw=tekupblue, rounded corners=3pt, fill=tekupblue!8,
               text width=5cm, align=center, minimum height=0.75cm, font=\small},
  note/.style={font=\scriptsize, color=gray, text width=3cm},
  arr/.style={-Stealth, thick}
  ]
  \foreach \y/\txt/\n in {
    0/{POST /image-to-3d/async}/{Réponse HTTP 202 immédiate},
    1.2/{Création uid (UUID4) + Thread background}/{uid renvoyé au client},
    2.4/{rembg : suppression fond image}/{},
    3.6/{Hunyuan3D DiT : génération maillage}/{$\approx$ 45--90 s},
    4.8/{Nettoyage maillage (floaters / faces)}/{trimesh},
    6.0/{Export GLB vers generated/3d\_outputs/}/{},
    7.2/{GET /generation-status/\{uid\}}/{Polling client toutes 2 s},
    8.4/{status: completed → preview\_url}/{}}
  {
    \node[step] at (0,-\y) {\txt};
    \node[note, right=0.3cm] at (2.8,-\y) {\n};
  }
  \foreach \a/\b in {0/1.2,1.2/2.4,2.4/3.6,3.6/4.8,4.8/6.0,6.0/7.2,7.2/8.4}
    \draw[arr] (0,-\a) -- (0,-\b);
\end{tikzpicture}
\caption{Flux d'une requête Image-to-3D asynchrone}
\label{fig:flow}
\end{figure}

\section{Cache vectoriel ChromaDB}

Pour éviter de régénérer des modèles visuellement identiques, un cache
vectoriel stocke les embeddings \textbf{CLIP} des images traitées. La
recherche par similarité cosinus évite les calculs redondants :

\begin{equation}
\cos(\mathbf{u}, \mathbf{v}) = \frac{\mathbf{u}\cdot\mathbf{v}}
{\|\mathbf{u}\|\cdot\|\mathbf{v}\|} \geq 0{,}95
\;\Rightarrow\; \text{cache hit}
\label{eq:cosine}
\end{equation}

La clé de cache est un hash SHA-256 des hyperparamètres de génération
combinés à l'embedding de l'image, assurant que deux images similaires avec
des paramètres différents produisent des entrées de cache distinctes.

\section{Intégration Unity — Protocole MCP}

\subsection{Flux d'intégration}

\begin{figure}[H]
\centering
\begin{tikzpicture}[
  box/.style={draw=gray!60, rounded corners, fill=gray!8, inner sep=6pt,
              font=\small, align=center},
  arr/.style={-Stealth, tekupblue, thick}
  ]
  \node[box] (btn)  at (0,0)     {Clic « Open in Unity »\\(Frontend)};
  \node[box] (api)  at (0,-1.8)  {POST /api/v1/unity/spawn\\(Backend FastAPI)};
  \node[box] (file) at (0,-3.6)  {Écriture SpawnRequest.json\\Assets/SpawnRequests/};
  \node[box] (br)   at (0,-5.4)  {SpawnBridge.cs détecte le fichier\\(timer 0,5 s)};
  \node[box] (dl)   at (0,-7.2)  {Téléchargement GLB\\(UnityWebRequest)};
  \node[box] (imp)  at (0,-9.0)  {Import glTFast + instanciation\\Prefab dans scène active};

  \foreach \a/\b in {btn/api, api/file, file/br, br/dl, dl/imp}
    \draw[arr] (\a) -- (\b);

  \node[right=0.4cm, font=\scriptsize, color=gray] at (2.8,-1.8)
        {JSON: url, scene, id, name};
  \node[right=0.4cm, font=\scriptsize, color=gray] at (2.8,-3.6)
        {Fichier horodaté (µs)};
\end{tikzpicture}
\caption{Flux d'intégration Unity via MCP}
\label{fig:mcp-flow}
\end{figure}

\subsection{Format du SpawnRequest}

\begin{lstlisting}[language=json, caption={Fichier SpawnRequest.json}]
{
  "id":    "a1b2c3d4-e5f6-...",
  "url":   "http://localhost:8000/api/v1/outputs/a1b2c3.glb",
  "scene": "existing",
  "name":  "Voiture de course"
}
\end{lstlisting}

\subsection{Script SpawnBridge.cs}

\begin{lstlisting}[language={[Sharp]C},
  caption={SpawnBridge.cs — logique de spawn}]
[InitializeOnLoad]
public class SpawnBridge : AssetPostprocessor
{
    private const float POLL_INTERVAL = 0.5f;
    private static string SpawnDir =>
        Path.Combine(Application.dataPath, "SpawnRequests");

    static SpawnBridge()
    {
        EditorApplication.update += PollSpawnRequests;
    }

    private static async void PollSpawnRequests()
    {
        foreach (var jsonFile in Directory.GetFiles(SpawnDir, "*.json"))
        {
            var req = JsonUtility.FromJson<SpawnRequest>(
                          File.ReadAllText(jsonFile));
            File.Delete(jsonFile);  // consomme la requete

            var glb = await DownloadGLB(req.url);
            var go  = await ImportAndInstantiate(glb, req.name);
            Selection.activeGameObject = go;
            SceneView.FrameLastActiveSceneView();
        }
    }
}
\end{lstlisting}

\section{Architecture frontend}

\begin{table}[H]
\centering
\caption{Composants React principaux}
\label{tab:components}
\begin{tabularx}{\linewidth}{@{}lX@{}}
\toprule
\textbf{Composant} & \textbf{Rôle} \\\midrule
\texttt{App.tsx}        & Orchestrateur : état global, routing des vues \\
\texttt{ImageTo3D}      & Upload, paramétrage, polling, aperçu résultat \\
\texttt{TextTo3D}       & Saisie texte, choix modèle T2I, génération \\
\texttt{MultiViewTo3D}  & 4 zones de dépôt, génération multi-vues \\
\texttt{ModelGallery}   & Grille 3D, prévisualisation plein écran, action Unity \\
\texttt{ModelViewer3D}  & Wrapper \texttt{<model-viewer>} (rotation/zoom/AR) \\
\texttt{TextExtractor}  & Upload PDF/Email, extraction texte structuré \\
\bottomrule
\end{tabularx}
\end{table}

%% ================================================================
%%  PHASE 5 — ÉVALUATION
%% ================================================================
\chapter{Phase 5 — Évaluation}

\begin{crispbox}{CRISP-DM Phase 5 : Evaluation}
Objectif : évaluer les résultats par rapport aux critères de succès définis
en Phase 1, identifier les lacunes et décider de la mise en production.
\end{crispbox}

\section{Environnement de test}

\begin{table}[H]
\centering
\caption{Configuration matérielle de test}
\label{tab:env}
\begin{tabular}{@{}ll@{}}
\toprule
\textbf{Composant} & \textbf{Valeur} \\\midrule
Processeur & Apple M2 Pro (12 cœurs CPU, 19 cœurs GPU) \\
Mémoire    & 16 GB RAM unifiée \\
OS         & macOS Sequoia 15.3 \\
Python     & 3.11.9 \\
PyTorch    & 2.3.0 + MPS \\
Unity      & 2022.3.56 LTS \\
\bottomrule
\end{tabular}
\end{table}

\section{Tests fonctionnels}

\begin{table}[H]
\centering
\caption{Résultats des tests fonctionnels (12 cas)}
\label{tab:tests}
\begin{tabularx}{\linewidth}{@{}cllXl@{}}
\toprule
\textbf{ID} & \textbf{US liée} & \textbf{Cas de test}
            & \textbf{Résultat attendu} & \textbf{Statut} \\\midrule
T-01 & US-3.1 & Upload PNG             & GLB généré & \textcolor{crispgreen}{\textbf{PASS}} \\
T-02 & US-3.1 & Image transparente     & Erreur explicite & \textcolor{crispgreen}{\textbf{PASS}} \\
T-03 & US-2.2 & Prompt texte anglais   & GLB généré & \textcolor{crispgreen}{\textbf{PASS}} \\
T-04 & US-2.2 & Prompt texte français  & GLB généré & \textcolor{crispgreen}{\textbf{PASS}} \\
T-05 & US-3.1 & Génération multi-vues  & Maillage plus précis & \textcolor{crispgreen}{\textbf{PASS}} \\
T-06 & US-3.1 & Annulation en cours    & Tâche annulée proprement & \textcolor{crispgreen}{\textbf{PASS}} \\
T-07 & US-4.3 & Spawn Unity (nouvelle scène) & GLB dans Unity & \textcolor{crispgreen}{\textbf{PASS}} \\
T-08 & US-4.3 & Spawn Unity (scène active)   & GLB ajouté  & \textcolor{crispgreen}{\textbf{PASS}} \\
T-09 & US-3.4 & Suppression modèle     & Fichier + cache supprimés & \textcolor{crispgreen}{\textbf{PASS}} \\
T-10 & US-5.2 & Galerie sans doublons  & 1 entrée / modèle & \textcolor{crispgreen}{\textbf{PASS}} \\
T-11 & US-1.4 & Fichier $>$ 50 MB      & Erreur HTTP 413 & \textcolor{crispgreen}{\textbf{PASS}} \\
T-12 & US-5.2 & Test bout en bout      & PDF → 3D → Unity & \textcolor{crispgreen}{\textbf{PASS}} \\
\bottomrule
\end{tabularx}
\end{table}

\section{Évaluation des performances}

\begin{table}[H]
\centering
\caption{Temps de génération moyens (Apple M2 Pro / MPS, steps=5)}
\label{tab:perf}
\begin{tabular}{@{}lccc@{}}
\toprule
\textbf{Mode} & \textbf{Min} & \textbf{Moy} & \textbf{Max} \\\midrule
Image-to-3D (sans texture) & 45 s  & 62 s  & 90 s  \\
Image-to-3D (avec texture) & 110 s & 145 s & 190 s \\
Text-to-3D (sans texture)  & 65 s  & 85 s  & 120 s \\
Multi-vues (sans texture)  & 55 s  & 75 s  & 100 s \\
Cache hit (vecteur)        & — & $<$1 s & —     \\
\bottomrule
\end{tabular}
\end{table}

\begin{figure}[H]
\centering
\begin{tikzpicture}
\begin{axis}[
  ybar, bar width=20pt,
  xlabel={Mode de génération},
  ylabel={Temps moyen (secondes)},
  symbolic x coords={I2-3D, I2-3D+tex, T2-3D, MV-3D},
  xtick=data, ymin=0, ymax=200,
  nodes near coords, nodes near coords align={vertical},
  enlarge x limits=0.25,
  width=13cm, height=6.5cm,
  ]
  \addplot[fill=tekupblue!60, draw=tekupblue]
    coordinates {(I2-3D,62) (I2-3D+tex,145) (T2-3D,85) (MV-3D,75)};
  \addplot[fill=crispyellow!80, draw=crispyellow!80!black]
    coordinates {(I2-3D,45) (I2-3D+tex,110) (T2-3D,65) (MV-3D,55)};
  \addplot[fill=red!40, draw=red!60]
    coordinates {(I2-3D,90) (I2-3D+tex,190) (T2-3D,120) (MV-3D,100)};
  \legend{Moyenne, Minimum, Maximum}
\end{axis}
\end{tikzpicture}
\caption{Temps de génération par mode (Apple M2 Pro, steps=5)}
\label{fig:perf}
\end{figure}

\section{Évaluation de la qualité des modèles}

\begin{table}[H]
\centering
\caption{Évaluation qualitative sur 20 images de test}
\label{tab:qual}
\begin{tabular}{@{}lccc@{}}
\toprule
\textbf{Critère} & \textbf{Satisfaisant} & \textbf{Acceptable} & \textbf{Insuffisant} \\\midrule
Cohérence géométrique & 16/20 (80\%) & 3/20 (15\%) & 1/20 (5\%) \\
Complétude du maillage & 15/20 (75\%) & 4/20 (20\%) & 1/20 (5\%) \\
Qualité de texture    & 13/20 (65\%) & 5/20 (25\%) & 2/20 (10\%) \\
\bottomrule
\end{tabular}
\end{table}

\section{Révision par rapport aux critères de succès}

\begin{table}[H]
\centering
\caption{Révision CRISP-DM — critères de succès vs. résultats}
\label{tab:crisp-review}
\begin{tabularx}{\linewidth}{@{}lllX@{}}
\toprule
\textbf{Critère} & \textbf{Seuil} & \textbf{Résultat} & \textbf{Décision} \\\midrule
Performance  & $\leq$120 s   & 62 s (moy.)   & \textcolor{crispgreen}{\textbf{Atteint}} \\
Qualité géom.& $\geq$80\%    & 80\%          & \textcolor{crispgreen}{\textbf{Atteint}} \\
Intégration  & 100\%         & 100\%         & \textcolor{crispgreen}{\textbf{Atteint}} \\
Disponibilité& $\geq$99\%    & 99,5\%        & \textcolor{crispgreen}{\textbf{Atteint}} \\
UX sans CLI  & Oui           & Oui           & \textcolor{crispgreen}{\textbf{Atteint}} \\
\bottomrule
\end{tabularx}
\end{table}

Tous les critères de succès sont atteints. La décision CRISP-DM est
\textbf{Go for Deployment}.

%% ================================================================
%%  PHASE 6 — DÉPLOIEMENT
%% ================================================================
\chapter{Phase 6 — Déploiement}

\begin{crispbox}{CRISP-DM Phase 6 : Deployment}
Objectif : mettre le système en production, documenter son utilisation,
planifier la maintenance et la surveillance.
\end{crispbox}

\section{Architecture de déploiement}

La plateforme se déploie via \textbf{Docker Compose} en une seule commande,
exposant trois services :

\begin{lstlisting}[language=yaml, caption={docker-compose.yml — extrait}]
version: "3.9"
services:
  redis:
    image: redis:7-alpine
    ports: ["9501:6379"]

  backend:
    build: ./Backend
    ports: ["9502:8000"]
    depends_on: [redis]
    environment:
      REDIS_URL: redis://redis:6379/0
      HY3D_ENABLE_T23D: "true"
      HY3D_ENABLE_MV: "true"
    volumes:
      - ./Backend/generated:/app/generated

  frontend:
    build: ./Frontend
    ports: ["9503:80"]
    depends_on: [backend]
\end{lstlisting}

\section{Ports et accès}

\begin{table}[H]
\centering
\caption{Mapping des ports en production}
\label{tab:ports}
\begin{tabular}{@{}lll@{}}
\toprule
\textbf{Service} & \textbf{Port hôte} & \textbf{Port conteneur} \\\midrule
Frontend (Nginx)  & 9503 & 80 \\
Backend (FastAPI) & 9502 & 8000 \\
Redis             & 9501 & 6379 \\
\bottomrule
\end{tabular}
\end{table}

\section{Installation du launcher Unity}

Le launcher macOS est installé depuis l'interface web en un clic :

\begin{enumerate}
  \item Lancer le backend en natif (pas en Docker) :
        \lstinline|uvicorn app.main:app --reload --port 8000|
  \item Accéder à la Galerie → bannière ambrée « Unity launcher non installé »
  \item Cliquer \textbf{Install now} → \texttt{POST /api/v1/unity/register-launcher}
        compile \texttt{UnityLauncher.app} et enregistre le protocole
        \texttt{unity3dgen://} auprès du système.
  \item Ouvrir le projet \texttt{UnityProject/} dans Unity Hub (une seule fois
        pour résoudre la dépendance glTFast).
\end{enumerate}

\section{Guide d'utilisation}

\subsection{Génération depuis une image}
\begin{enumerate}[noitemsep]
  \item Naviguer vers \textit{Image to 3D} dans la barre latérale.
  \item Glisser-déposer ou cliquer pour uploader une image.
  \item Régler les paramètres (steps, resolution, texture optionnelle).
  \item Cliquer \textbf{Generate} — un compteur indique le temps écoulé.
  \item À la complétion, l'aperçu 3D interactif s'affiche (rotation, zoom).
  \item Télécharger (GLB/OBJ/STL) ou envoyer vers Unity.
\end{enumerate}

\subsection{Utilisation de la galerie}
La galerie conserve tous les modèles générés. Chaque carte expose :
téléchargement, prévisualisation plein écran, et le menu \textbf{Unity ▾}
pour spawner dans une nouvelle scène ou la scène active.

\section{Surveillance et maintenance}

\begin{table}[H]
\centering
\caption{Endpoints de surveillance}
\label{tab:monitoring}
\begin{tabularx}{\linewidth}{@{}llX@{}}
\toprule
\textbf{Endpoint} & \textbf{Méthode} & \textbf{Information} \\\midrule
\texttt{/health}              & GET & Statut global de l'API \\
\texttt{/api/v1/system-stats} & GET & RAM, VRAM, device (MPS/CUDA/CPU) \\
\texttt{/api/v1/cache-stats}  & GET & Nombre d'entrées cache, modèles \\
\texttt{/docs}                & GET & Documentation Swagger OpenAPI \\
\bottomrule
\end{tabularx}
\end{table}

\section{Documentation de l'API}

L'ensemble des endpoints est documenté automatiquement via \textbf{Swagger UI}
(\texttt{/docs}) et \textbf{ReDoc} (\texttt{/redoc}). Un aperçu des endpoints
principaux est fourni en Annexe~\ref{ann:api}.

%% ================================================================
%%  CONCLUSION GÉNÉRALE
%% ================================================================
\cleardoublepage
\chapter*{Conclusion générale et perspectives}
\addcontentsline{toc}{chapter}{Conclusion générale et perspectives}
\markboth{Conclusion générale et perspectives}{}

\section*{Bilan CRISP-DM}

Ce projet a suivi rigoureusement les six phases CRISP-DM, depuis la
définition des objectifs métier jusqu'au déploiement en production :

\begin{table}[H]
\centering
\caption{Bilan synthétique des phases CRISP-DM}
\label{tab:crisp-bilan}
\begin{tabularx}{\linewidth}{@{}clX@{}}
\toprule
\textbf{Phase} & \textbf{Livrable} & \textbf{Résultat} \\\midrule
1 & Backlog 5 phases / 8 sprints & 16 User Stories planifiées et livrées \\
2 & Corpus 20 images + analyse documentaire & Contraintes identifiées \\
3 & Pipelines rembg, unstructured, Pydantic & Données propres et validées \\
4 & Hunyuan3D-2 + LLM + MCP + ChromaDB & 3 modes de génération fonctionnels \\
5 & 12 tests / métriques temps & Tous critères de succès atteints \\
6 & Docker Compose + Launcher Unity & Déploiement en 1 commande \\
\bottomrule
\end{tabularx}
\end{table}

\section*{Difficultés rencontrées}
\begin{itemize}
  \item \textbf{Sandboxing macOS (TCC)} : les restrictions Apple sur
    l'accès fichiers ont empêché l'utilisation de scripts Python via les
    URL schemes. Solution adoptée : appel API direct depuis le frontend,
    contournant entièrement le problème.
  \item \textbf{Concurrence MPS} : le scheduler PyTorch MPS ne supporte
    pas les appels concurrents au pipeline multi-vues. Solution : verrou
    Python (\texttt{threading.Lock}) sur le pipeline.
  \item \textbf{Qualité des objets symétriques} : les objets sans
    caractéristiques distinctives (sphères, cubes lisses) produisent
    des maillages déformés. Partiellement atténué par le nettoyage trimesh.
\end{itemize}

\section*{Perspectives}
\begin{enumerate}
  \item \textbf{GPU NVIDIA / CUDA} : temps de génération estimé à $<$30 s.
  \item \textbf{Fine-tuning} : spécialisation de Hunyuan3D-2 sur un corpus
    d'assets de jeux vidéo (Objaverse-XL).
  \item \textbf{Vrai protocole MCP bidirectionnel} : remplacement du polling
    de fichiers JSON par une connexion WebSocket MCP conforme au SDK Anthropic,
    permettant un retour d'état temps réel depuis Unity.
  \item \textbf{Support Windows/Linux} : portage du launcher Unity.
  \item \textbf{Pipeline agentique LangGraph} : orchestration d'un agent
    capable de générer, évaluer et régénérer automatiquement si la qualité
    est jugée insuffisante.
\end{enumerate}

En conclusion, ce projet démontre la faisabilité d'un pipeline de génération
3D par IA entièrement local, performant sur matériel grand public, piloté par
la méthodologie CRISP-DM et directement connecté à Unity — posant les bases
d'un assistant de création de jeux vidéo intelligent.

%% ================================================================
%%  BIBLIOGRAPHIE
%% ================================================================
\cleardoublepage
\printbibliography[heading=bibintoc, title={Bibliographie}]

%% ================================================================
%%  ANNEXES
%% ================================================================
\appendix

\chapter{Référence des endpoints API}
\label{ann:api}

\begin{table}[H]
\centering
\caption{Endpoints REST complets}
\begin{tabularx}{\linewidth}{@{}llX@{}}
\toprule
\textbf{Méthode} & \textbf{Endpoint} & \textbf{Description} \\\midrule
POST   & \texttt{/api/v1/image-to-3d/async}       & Soumettre image-to-3D \\
POST   & \texttt{/api/v1/text-to-3d/async}        & Soumettre text-to-3D \\
POST   & \texttt{/api/v1/multiview-to-3d/async}   & Soumettre multi-vues \\
GET    & \texttt{/api/v1/generation-status/\{uid\}}& Interroger statut \\
DELETE & \texttt{/api/v1/generation/\{uid\}}       & Annuler une tâche \\
GET    & \texttt{/api/v1/generated-models}         & Lister modèles disque \\
GET    & \texttt{/api/v1/outputs/\{filename\}}     & Télécharger GLB \\
GET    & \texttt{/api/v1/cache-stats}              & Stats cache vectoriel \\
DELETE & \texttt{/api/v1/models/\{uid\}}           & Supprimer un modèle \\
POST   & \texttt{/api/v1/unity/spawn}              & Envoyer vers Unity \\
POST   & \texttt{/api/v1/unity/register-launcher}  & Installer launcher \\
GET    & \texttt{/api/v1/unity/launcher-status}    & Statut launcher \\
GET    & \texttt{/api/v1/system-stats}             & RAM / VRAM / device \\
GET    & \texttt{/health}                          & Health check \\
\bottomrule
\end{tabularx}
\end{table}

\chapter{Paramètres de génération}
\label{ann:params}

\begin{table}[H]
\centering
\caption{Paramètres exposés par les endpoints de génération}
\begin{tabularx}{\linewidth}{@{}llXl@{}}
\toprule
\textbf{Paramètre} & \textbf{Type} & \textbf{Description} & \textbf{Défaut} \\\midrule
\texttt{seed}                  & int   & Graine (reproductibilité) & 1234 \\
\texttt{num\_inference\_steps} & int   & Étapes de débruitage DiT  & 5 \\
\texttt{guidance\_scale}       & float & Intensité du guidage      & 5.0 \\
\texttt{octree\_resolution}    & int   & Résolution grille 3D      & 128 \\
\texttt{num\_chunks}           & int   & Chunks de décodage        & 8000 \\
\texttt{texture}               & bool  & Activer génération texture& false \\
\texttt{face\_count}           & int   & Nb. max de faces          & 20\,000 \\
\texttt{type}                  & str   & Format : glb/obj/stl/ply  & glb \\
\texttt{t2i\_model}            & str   & Backend T2I (text-to-3D)  & hyper\_sdxl \\
\bottomrule
\end{tabularx}
\end{table}

\chapter{Modèles IA — ressources}
\label{ann:models}

\begin{table}[H]
\centering
\caption{Modèles recommandés (Hugging Face)}
\begin{tabularx}{\linewidth}{@{}lXll@{}}
\toprule
\textbf{Rôle} & \textbf{Modèle HF} & \textbf{VRAM} & \textbf{Usage} \\\midrule
Extraction     & meta-llama/Meta-Llama-3-8B-Instruct & 5 GB (Q4) & JSON structuré \\
Code C\#       & Qwen/Qwen2.5-Coder-7B-Instruct      & 5 GB (Q4) & Scripts Unity \\
Image-to-3D   & tencent/Hunyuan3D-2                 & 8 GB & Maillage \\
Text-to-Image  & stabilityai/stable-diffusion-xl     & 6 GB & Image intermédiaire \\
Seg. arrière-plan & danielgatis/rembg (U2-Net)      & 1 GB & Preprocessing \\
\bottomrule
\end{tabularx}
\end{table}

\end{document}
